% Sample-size sensitivity of NMI under a null where all classes are i.i.d.
% Source of record: experiments/analysis/clustering/nmi_null_simulation.py (seed 20260923)
%   -> experiments/results/clustering/null/nmi_null_simulation.{json,md}; T=1000 per row.
% Supports the clustering cross-check (Section~\ref{sec:eval-clustering}): explains why
% small-sample NMI can be elevated even with no class-distinguishing signal.
\section{Sample-Size Sensitivity of NMI}
\label{app:nmi-sensitivity}

The clustering cross-check reads ARI directly but debiases the NMI of the balanced PI pairs (Section~\ref{sec:eval-method}).
To examine that choice, we quantify how \emph{finite-sample effects} alone can inflate NMI
even when there is \emph{no} discriminative signal, using an idealized null in which all
classes are statistically indistinguishable.

We sample every class from the \emph{same} distribution, $z_{i}\sim\mathcal{N}(0,I_{d})$ for all classes, where $I_d$ is the $d$-dimensional identity covariance and $d=2$ matches the two gas features the clustering uses.
Since no class carries class-specific information, a clustering algorithm should not recover class membership,
and both NMI and ARI should concentrate around $0$ as the sample size grows.

For each configuration we run $T=1000$ independent trials.
In each trial we
(i)~generate $G$ classes with $n$ i.i.d.\ samples per class from $\mathcal{N}(0,I_d)$, for $n_{\mathrm{tot}}=Gn$ samples in total, assigning each sample a class label by construction order (these labels are independent of the sample values, since all classes share one distribution),
(ii)~standardize the pooled data and apply K-Means with $K=G$ and ten initializations, as in the evaluation, and
(iii)~score the predicted assignments against these assigned class labels with NMI and ARI.
We vary $n$ while holding $G$ fixed, and report the trial means $\overline{\mathrm{NMI}}$ and $\overline{\mathrm{ARI}}$.
We consider a \textbf{PI-like} binary structure ($G=2$, $n\in\{3,10,50,100\}$), a five-class structure ($G=5$, $n\in\{3,10,100\}$), and a ten-class structure ($G=10$, $n\in\{3,10,77\}$) together with the observed ICD-10 chapter sizes ($49$--$135$ executions per chapter, $n_{\mathrm{tot}}=770$).

Table~\ref{tab:nmi-sample-size-sensitivity} shows that, although all classes are indistinguishable by construction, \textbf{NMI is noticeably positive at very small $n$ but collapses toward zero as $n$ grows}, whereas \textbf{ARI stays near zero throughout}.
For the binary structure, $\overline{\mathrm{NMI}}$ falls from $0.2074$ at $n=3$ ($n_{\mathrm{tot}}=6$) to $0.0420$ at $n=10$, and to $0.0072$/$0.0034$ at $n=50$/$100$.
The same trend holds with more classes, at a higher level for every $n$: with five classes $\overline{\mathrm{NMI}}$ goes from $0.4152$ ($n=3$) to $0.0102$ ($n=100$), and with ten classes from $0.5346$ ($n=3$) to $0.0239$ ($n=77$).
For the observed ICD-10 chapter sizes, a separate null run that averages each trial over $20$ K-Means seeds, as the evaluation does, gives a null NMI with mean $0.0246$ and $95$th percentile $0.0302$; the observed ICD-10 NMI of $0.0281$ is reached in $145$ of $1000$ null trials ($p=0.15$).

\begin{table}[htbp]
	\centering
	\caption{Sample-size sensitivity of NMI and ARI under the null hypothesis where all
		classes are i.i.d.\ sampled from $\mathcal{N}(0,I_d)$.}
	\label{tab:nmi-sample-size-sensitivity}
	\begin{tabular*}{0.54\linewidth}{@{\extracolsep{\fill}}rrrr@{}}
		\toprule
		$n$ (per class) & $n_{\mathrm{tot}}$ & $\overline{\mathrm{NMI}}$ & $\overline{\mathrm{ARI}}$ \\
		\midrule
		\multicolumn{4}{@{}l}{\textit{Two classes (as in a PI pair)}} \\
		$3$       & $6$   & $0.2074$ & $0.0127$  \\
		$10$      & $20$  & $0.0420$ & $0.0008$  \\
		$50$      & $100$ & $0.0072$ & $-0.0003$ \\
		$100$     & $200$ & $0.0034$ & $-0.0004$ \\
		\addlinespace
		\multicolumn{4}{@{}l}{\textit{Five classes}} \\
		$3$       & $15$  & $0.4152$ & $-0.0034$ \\
		$10$      & $50$  & $0.1226$ & $0.0003$  \\
		$100$     & $500$ & $0.0102$ & $-0.0000$ \\
		\addlinespace
		\multicolumn{4}{@{}l}{\textit{Ten classes (as the DUO and ICD-10 targets)}} \\
		$3$       & $30$  & $0.5346$ & $-0.0009$ \\
		$10$      & $100$ & $0.2087$ & $0.0003$  \\
		$77$      & $770$ & $0.0239$ & $-0.0001$ \\
		\addlinespace
		\multicolumn{4}{@{}l}{\textit{Ten classes, observed ICD-10 chapter sizes}} \\
		$49$--$135$ & $770$ & $0.0244$ & $-0.0001$ \\
		\bottomrule
	\end{tabular*}

	\begin{tablenote}
		$T=1000$ trials per row, one K-Means fit with ten initializations per trial; $n_{\mathrm{tot}}$ is the total sample size.
		For the observed ICD-10 chapter sizes, the text also reports a separate run that averages each trial over $20$ K-Means seeds, as the evaluation does.
	\end{tablenote}
\end{table}
